\chapter{Einleitung}
\label{ch:einleitung}

% Motivation: warum ist Offline-Zugang zu Wikipedia relevant (Ausfaelle,
% fehlende Internetanbindung, Katastrophenszenarien -- siehe Aufgabenstellung).
% Kurz die konkrete Aufgabenstellung zusammenfassen: Auswahl einer Teilmenge
% + dichte Vektorindexierung + RAG auf einem Raspberry Pi 5 (16 GB RAM, 1 TB SSD).

\section{Motivation und Problemstellung}

Wikipedia stellt einen umfangreichen Wissensbestand bereit, dessen Nutzung
üblicherweise eine funktionierende Internetverbindung voraussetzt. In
Situationen, in denen kein direkter Internetzugang verfügbar ist, beispielsweise
aufgrund technischer Störungen oder fehlender Netzanbindung, ist der Zugriff auf
diesen Wissensbestand daher nur eingeschränkt möglich. Ziel des Projekts WikiPod
ist es, zu untersuchen, inwieweit ein Teil der englischsprachigen Wikipedia
lokal bereitgestellt und ohne direkte Internetverbindung nutzbar gemacht werden
kann.

Eine zentrale Randbedingung stellt dabei die vorgesehene Zielhardware dar. Das
System soll auf einem Raspberry Pi~5 mit 16~GB Arbeitsspeicher und einer
1-TB-SSD betrieben werden. Aufgrund dieser begrenzten Ressourcen kann nicht
davon ausgegangen werden, dass sämtliche Daten und Verarbeitungsschritte ohne
weitere Einschränkungen auf die Zielhardware übertragen werden können. Es muss
daher zunächst eine geeignete Teilmenge der englischsprachigen Wikipedia
ausgewählt werden, die ein vorgegebenes Speicherbudget einhält.

Neben der lokalen Bereitstellung soll der ausgewählte Wissensbestand auch für
natürlichsprachliche Anfragen nutzbar sein. Hierzu werden die ausgewählten
Wikipedia-Inhalte als dichte Vektoren in OpenSearch indexiert und relevante
Textpassagen zu einer Anfrage abgerufen. Diese bilden anschließend den Kontext
für ein lokal betriebenes Sprachmodell im Rahmen eines
Retrieval-Augmented-Generation-Ansatzes. Neben der funktionalen Umsetzung stellt
sich damit insbesondere die Frage, wie effektiv das Retrieval arbeitet und
inwieweit sich die Verarbeitung unter den gegebenen Hardwarebeschränkungen
praktikabel durchführen lässt.

\section{Ziele der Projektarbeit}

Ausgehend von dieser Problemstellung verfolgt die Projektarbeit fünf zentrale
Ziele. Zunächst soll eine lokale Kopie ausgewählter Inhalte der
englischsprachigen Wikipedia bereitgestellt werden. Da die verfügbare
Speicherkapazität begrenzt ist, soll hierfür ein Verfahren entwickelt werden,
das eine geeignete Teilmenge der Wikipedia unter Einhaltung eines vorgegebenen
Speicherbudgets auswählt.

Die ausgewählten Inhalte sollen anschließend als dichte Vektoren in OpenSearch
indexiert werden, sodass zu einer natürlichsprachlichen Anfrage relevante
Textpassagen semantisch abgerufen werden können. Auf dieser Grundlage soll ein
für die eingeschränkte Zielhardware geeignetes Sprachmodell ausgewählt und mit
den abgerufenen Textpassagen zur lokalen Antwortgenerierung eingesetzt werden.

Abschließend soll die entwickelte Lösung sowohl hinsichtlich ihrer Effektivität
als auch ihrer Effizienz evaluiert werden. Dabei werden insbesondere die
Qualität des Retrievals sowie Laufzeit und Ressourcenbedarf auf der
vorgesehenen Hardware betrachtet. Die Projektarbeit untersucht damit die
folgenden fünf Schwerpunkte:

\begin{enumerate}
    \item Bereitstellung einer lokalen Kopie der englischsprachigen Wikipedia,
    \item Auswahl einer Teilmenge unter Berücksichtigung eines Speicherbudgets,
    \item Indexierung der ausgewählten Inhalte als dichte Vektoren mit OpenSearch,
    \item Auswahl und Einsatz eines für die Zielhardware geeigneten Sprachmodells,
    \item Evaluation der Gesamtlösung hinsichtlich Effektivität und Effizienz.
\end{enumerate}

\section{Aufbau des Berichts}

Der weitere Bericht folgt den wesentlichen Verarbeitungsschritten und
Fragestellungen des Projekts. Kapitel~\ref{ch:architektur} gibt zunächst einen
Überblick über die Systemarchitektur, die eingesetzten Technologien und den
Aufbau der Verarbeitungspipeline.

Anschließend beschreibt Kapitel~\ref{ch:dokumentenauswahl} die Auswahl einer
Teilmenge der Wikipedia unter Berücksichtigung des verfügbaren Speicherbudgets.
Kapitel~\ref{ch:indexierung} behandelt die Aufbereitung der ausgewählten
Artikel, ihre Einbettung als dichte Vektoren sowie die Indexierung und das
Retrieval mit OpenSearch. Darauf aufbauend wird die lokale Antwortgenerierung
mithilfe eines Sprachmodells im Rahmen des RAG-Ansatzes erläutert.

Die Umsetzung unter den Ressourcenbeschränkungen des Raspberry Pi sowie die
dabei auftretenden technischen Herausforderungen werden anschließend im
Deployment-Kapitel betrachtet. Die Evaluation untersucht die entwickelte
Lösung hinsichtlich Retrieval-Qualität, Laufzeit und Ressourcenbedarf. Eine
Anleitung zur Inbetriebnahme fasst die für die Nutzung und Reproduktion
notwendigen Schritte zusammen.

Abschließend werden die getroffenen Designentscheidungen und die Grenzen der
Lösung diskutiert. Das Fazit fasst die Ergebnisse der Projektarbeit zusammen
und zeigt mögliche Ansatzpunkte für zukünftige Weiterentwicklungen auf.
